%% kapa: exam=1 round=27 year=2016 period=1 form=A subject=civil_law subject_ko=민법 seq=8 number=8
%% source: Q-Net 감정평가사 기출문제 (cst00309) artlSeq=5169402 file=extracted/1008/5169402/제27회 감정평가사 1차 문제/제27회 감정평가사 1교시문제(A형).pdf
%% conversion: kordoc 4.24.51 -> tools/gen1.py (auto); review=proofread-sonnet; flags=none
\begin{question}{8}
착오에 관한 설명으로 옳지 않은 것은? (다툼이 있으면 판례에 따름)
\begin{choices}
  \item 대리인에 의한 법률행위에서 착오의 유무는 대리인을 표준으로 판단한다.
  \item 착오의 존재 여부는 의사표시 당시를 기준으로 판단한다.
  \item 착오에 의한 취소의 의사표시는 반드시 명시적이어야 하는 것은 아니고, 취소자가 그 착오를 이유로 자신의 법률행위의 효력을 처음부터 배제하려고 한다는 의사가 드러나면 충분하다.
  \item 착오가 표의자의 중대한 과실로 인한 경우, 상대방이 표의자의 착오를 알고 이를 이용하였더라도, 표의자는 의사표시를 취소할 수 없다.
  \item 착오를 이유로 법률행위를 취소한 표의자가 상대방에게 불법행위책임을 지는 것은 아니다.
\end{choices}
\end{question}
\answer{4}
